% Generated draft: requires mathematical and source review.
\documentclass[11pt]{article}
\usepackage[margin=25mm]{geometry}
\usepackage{fontspec}
\setmainfont{Latin Modern Roman}
\usepackage{amsmath,amssymb,mathtools,unicode-math}
\setmathfont{Latin Modern Math}
\usepackage{enumitem,needspace,xurl,hyperref}
\setlength{\parindent}{0pt}
\setlength{\parskip}{5pt}
\setlength{\emergencystretch}{3em}
\newcommand{\N}{\mathbb{N}}
\newcommand{\Z}{\mathbb{Z}}
\newcommand{\Q}{\mathbb{Q}}
\newcommand{\R}{\mathbb{R}}
\newcommand{\C}{\mathbb{C}}
\newcommand{\F}{\mathbb{F}}
\newcommand{\A}{\mathbb{A}}
\newcommand{\PP}{\mathbb{P}}
\newcommand{\E}{\mathbb{E}}
\newcommand{\eps}{\varepsilon}
\newcommand{\dd}{\,\mathrm{d}}
\begin{document}
\section*{UnsolvedMath ID: 30001735}
\textbf{Source problem code:} catalogue entry
\subsection{Definitions and mathematical statement}
The supplied record concerns two orthogonal-group \(L\)-packets, denoted here by \(\Pi_\phi\) and \(\Pi_{\phi'}\), with Langlands parameters \(\phi\) and \(\phi'\). In the parametrizations referred to by the source, \(\pi(\phi,\chi)\) denotes the representation labeled by a component-group character \(\chi\), and similarly for \(\phi'\). The symbols \(\epsilon_{\phi,\phi'}\) and \(\epsilon'_{\phi,\phi'}\) denote the Gross--Prasad characters that the source describes as defined by epsilon factors. Their formulas and normalization conventions are not supplied.

Write \(m(\pi,\pi')\) schematically for the Gross--Prasad multiplicity intended in the source. The extracted conjecture asserts that, under the intended Langlands parametrizations, the unique pair of multiplicity one is
\[
\bigl(\pi(\phi,\epsilon_{\phi,\phi'}),\,
\pi(\phi',\epsilon'_{\phi,\phi'})\bigr).
\]
Equivalently, if \(\mathcal P_{\phi,\phi'}\) denotes the intended collection of admissible pairs from the two packets, the assertion has the schematic form
\[
\bigl\{(\pi,\pi')\in\mathcal P_{\phi,\phi'}:
 m(\pi,\pi')=1\bigr\}
=
\left\{\bigl(\pi(\phi,\epsilon_{\phi,\phi'}),\,
\pi(\phi',\epsilon'_{\phi,\phi'})\bigr)\right\}.
\]
Here \(\mathcal P_{\phi,\phi'}\) and \(m\) are placeholders for missing source context, not independently specified mathematical objects. In particular, this display does not additionally assert that every other multiplicity is zero.

The supplied extraction does not specify the base field, the precise orthogonal groups and their relation, the allowed parameters, the treatment of inner forms, the admissible-pair collection, the multiplicity space, or the normalizations of the parametrizations and epsilon factors. Consequently, it does not determine a fully quantified conjecture without further source context. The assertion above preserves the extracted target conditionally on those unresolved choices.
\subsection{Short English statement}
The extracted Gross--Prasad conjecture identifies the unique multiplicity-one pair in two orthogonal-group \(L\)-packets by the characters prescribed by epsilon factors. Essential hypotheses and normalization conventions are missing, so this is a schematic statement requiring human mathematical review. The supplied partially solved label is dated curation, not an independently verified conclusion.
\subsection{Sources}
\begin{itemize}
\item Supplied local UnsolvedMath registry record, permanent ID \(30001735\), exact source problem code \texttt{catalogue entry}, title ``Gross--Prasad Labels for Multiplicity-One Pairs,'' in the Oberwolfach Reports Open Problems collection. No individual proposer is identified in the supplied record.
\item \emph{Automorphic Forms: New Directions} (2011), Oberwolfach Reports, pp.\ 719--779, DOI \texttt{10.4171/owr/2011/14}. Supplied source report and literature reference labeled ``original OWR report'': \url{https://doi.org/10.4171/owr/2011/14}. The supplied background describes Oberwolfach Reports as documenting mathematical workshops and their open-problem discussions.
\item Supplied literature reference labeled ``Gan--Gross--Prasad'': \url{https://doi.org/10.24033/ast.912}. No more detailed bibliographic information is supplied for this reference.
\item Supplied literature reference labeled ``tempered special-orthogonal case'': \url{https://arxiv.org/abs/2009.13947}. No more detailed bibliographic information is supplied for this reference.
\item Literature assessment dated exactly \texttt{2026-08-21}: the supplied status is ``partially solved.'' The supplied assessment states that the component-group parametrization and epsilon labeling are established in substantial classical-group cases, while the report's broad Langlands-compatibility assertion is not a single fully settled theorem in the extracted scope. This is supplied dated curation derived from the OWR material, not independently verified here; time-sensitive status and cited claims require independent verification.
\item The supplied registry labels the statement \texttt{corrected\_verified} and reports that its mathematical content was checked against the official Oberwolfach report, with notation, context, and question form normalized without changing the source-level claim. This verification report is supplied provenance and has not been independently checked here.
\end{itemize}
\end{document}
